%
% Exponentialdarstellung der trigonometrischen Funktionen
%
\subsection{Exponentialdarstellung der trigonometrischen Funktionen
\label{buch:intalg:elementar:subsection:trigo}}
Die trigonometrischen Funktionen wurden ursprünglich für die Zwecke
der Beschreibung der Himmelskugel entwickelt und in der ebenen
Geometrie bei der Berechnung von Dreiecken angewendet.
Mit der eulerschen Formel $e^{it}=\cos t + i\sin t$ wird jedoch 
\index{eulersche Formel}%
auch eine Beziehung zur komplexen Exponentialfunktion offengelegt.
Sie ermöglicht, die trigonometrischen Funktionen und ihre
Umkehrfunktionen ausschliesslich über die Exponentialfunktion und
ihre Umkehrfunktion, den Logarithmus, zu beschreiben.

%
% Hyperbolische Funktionen
%
\subsubsection{Hyperbolische Funktionen}
Bei den hyperbolischen Funktionen ist die Situation viel einfacher,
da diese Funktionen bereits durch die reelle Exponentialfunktion
definiert sind.

\begin{definition}[Hyperbolische Funktionen]
\label{buch:intalg:elementar:trigo:def:hyperbolisch}
Die \emph{hyperbolischen Funktionen} sind durch
\index{hyperbolische Funktion}%
\begin{align*}
\sinh x &= \frac{e^x-e^{-x}}2,
&
\cosh x &= \frac{e^x+e^{-x}}2,
&
\tanh x &= \frac{e^x-e^{-x}}{e^x+e^{-x}},
&
\coth x &= \frac{e^x+e^{-x}}{e^x-e^{-x}}
\end{align*}
für beliebige Argumente $x\in \mathbb{C}$ definiert.
\index{sinh@$\sinh$}%
\index{cosh@$\cosh$}%
\index{tanh@$\tanh$}%
\index{coth@$\coth$}%
\end{definition}

Die hyperbolischen Funktionen sind besonderes nützlich für die Lösung
von Differentialgleichungen zweiter Ordnung, weil sie ähnlich einfache 
Anfangsbedingungen wie die Funktionen $\cos x$ und $\sin x$ haben.
Die Differentialgleichung $y''=y$ hat die Exponentialfunktionen 
$y_+(x)=e^x$ und $y_-(x)=e^{-x}$ als linear unabhängige Lösungen.
Um jedoch eine Lösung der Differentialgleichung zu den Anfangsbedingungen
$y(0)=y_0$ und $y'(0)=y_1$ zu finden, muss man eine Linearkombination
\[
y(x) = A y_+(x) + B y_-(x)
\]
finden, die die korrekten Anfangsbedingungen erfüllt. 
Dazu setzt  man die Linearkombination in die Anfangsbedingungen ein
und erhält das Gleichungssystem
\index{lineares Gleichungssystem}%
\[
\renewcommand{\arraycolsep}{2pt}
\left.
\begin{array}{rcrcrcr}
y(0)  &=& A y_+(0)  &+& B y_-(0)  &=& y_0 \\
y'(0) &=& A y_+'(0) &+& B y_-'(0) &=& y_1 \\
\end{array}
\quad
\right\}
\qquad
\Rightarrow
\qquad
\left\{
\quad
\begin{array}{rcrcl}
A &+& B &=& y_0 \\
A &-& B &=& y_1.
\end{array}
\right.
\]
Durch Addition bzw.~Subtraktion der beiden Gleichungen ist das 
Gleichungssystem einfach zu lösen, man findet die Koeffizienten
\[
A
=
\frac{y_0+y_1}2
\qquad\text{und}\qquad
B
=
\frac{y_0-y_1}2
\]
und daraus die Lösung
\begin{align*}
y(x)
=
Ay_+(x) + By_-(x)
&=
\frac{y_0+y_1}2 e^x
+
\frac{y_0-y_1}2 e^{-x}
\\
&=
y_0
\frac{e^x+e^{-x}}2
+
y_1
\frac{e^x-e^{-x}}2
=
y_0\cosh x + y_1 \sinh x
\end{align*}
der Differentialgleichung.
Die hyperbolischen Funktionen tauchen also genauso natürlich in der
Lösung der Differentialgleichung $y''=y$ auf, wie die trigonometrischen
Funktionen in der Lösung der Schwingungsdifferentialgleichung $y''=-y$
auftreten.
Der Grund dafür ist natürlich, dass die Funktionen die gleichen
Anfangsbedingungen
\[
\begin{aligned}
\cos  0 &= 1,& \cos'  0 &= 0,\\
\sin  0 &= 0,& \sin'  0 &= 1
\end{aligned}
\qquad\text{bzw.}\qquad
\begin{aligned}
\cosh 0 &= 1,& \cosh' 0 &= 0,\\
\sinh 0 &= 0,& \sinh' 0 &= 1
\end{aligned}
\]
haben.

%
% Die hyperbolischen Umkehrfunktionen
%
\subsubsection{Die hyperbolischen Umkehrfunktionen}
Die Logarithmusfunktion ist die Umkehrfunktion der Exponentialfunktion.
Mit ihr lassen sich auch die hyperbolischen Funktionen invertieren.

Um die Umkehrfunktion $\operatorname{arcosh}y$ zu finden, muss die
reelle Zahl $x$ gefunden werden, die $\cosh x = y$ gilt.
Aus der Definition folgt die Gleichung
\begin{align*}
\frac{e^x + e^{-x}}2 &= y.
\intertext{Da die Exponentialfunktion nicht verschwinden kann, verlieren
wir nichts, wenn wir mit $e^x$ multiplizieren. 
Dies ergibt}
\frac{(e^x)^2 + 1}2 &= y e^x
\intertext{oder als quadratische Gleichung für $e^x$ geschrieben:}
(e^x)^2 - 2y e^x + 1 &= 0.
\end{align*}
Die Lösung der quadratischen Gleichung ist
\begin{equation}
e^x
=
y\pm \!\sqrt{\smash{y^2}-1\raisebox{1pt}{\mathstrut}}
\qquad
\Rightarrow
\qquad
x
=
\log\Bigl(
y\pm \!\sqrt{\smash{y^2}-1\raisebox{1pt}{\mathstrut}}
\Bigr).
\label{buch:intalg:elementar:eqn:arcosh}
\end{equation}
Da $\cosh x$ eine gerade Funktion ist, gibt es zu fast allen
$y$-Werten zwei $x$-Werte.
Wir würden auch erwarten, dass die beiden Werte in
\eqref{buch:intalg:elementar:eqn:arcosh} entgegengesetzt sind.
Dazu müssen aber die Argumente des Logarithmus reziprok sein.
Tatsächlich folgt durch Rationalisieren
\begin{align*}
\frac{1}{y\pm\!\sqrt{\smash{y^2}-1\raisebox{1.5pt}{\mathstrut}}}
&=
\frac{
y\mp\!\sqrt{\smash{y^2}-1\raisebox{1.5pt}{\mathstrut}}
}{
(
y\pm\!\sqrt{\smash{y^2}-1\raisebox{1.5pt}{\mathstrut}}
)(
y\mp\!\sqrt{\smash{y^2}-1\raisebox{1.5pt}{\mathstrut}}
)
}
\\
&=
\frac{
y\mp\!\sqrt{\smash{y^2}-1\raisebox{1.5pt}{\mathstrut}}
}{
y^2-(y^2-1)
}
\\
&=
y\mp\!\sqrt{\smash{y^2}-1\raisebox{1.5pt}{\mathstrut}}.
\end{align*}
Die beiden Argumente des Logarithmus sind also tatsächlich
reziprok.

Auch für die hyperbolische Sinusfunktion kann mit der gleichen
Methode ein Ausdruck für die Umkehrfunktion gefunden werden.
Aus
\[
\sinh x = y
\qquad\Rightarrow\qquad
\frac{e^x - e^{-x}}2 = y
\]
folgt durch Multiplikation mit $2e^x$
\begin{align*}
(e^x)^2 -2ye^x - 1 &= 0
\notag
\intertext{mit der Lösung}
e^x &= y\pm\!\sqrt{\smash{y^2}+1\raisebox{1.5pt}{\mathstrut}}.
\label{buch:intalg:elementar:trig:eqn:arsinh}
\end{align*}
Da das negative Zeichen auf der rechten Seite zu einem negativen Wert 
führt, bleibt nur die Lösung
\[
x
=
\log\Bigl(y+\!\sqrt{\smash{y^2}+1\raisebox{1.5pt}{\mathstrut}}\Bigr)
.
\]

Für die hyperbolische Tangens- und Kotangensfunktion ist der
entsprechende Ansatz
\begin{align*}
\tanh x
=
\frac{e^x-e^{-x}}{e^x+e^{-x}}
&=
y
&
&\text{bzw.}&
\coth x
=
\frac{e^x+e^{-x}}{e^x-e^{-x}}
&=
y
\intertext{oder nach Multiplizieren mit dem Nenner}
e^x-e^{-x}
&=
ye^x+ye^{-x}
&
&&
e^x+e^{-x}
&=
ye^x-ye^{-x}
\intertext{und $e^x$}
(e^x)^2-1
&=
y(e^x)^2+y
&
&&
(e^x)^2+1
&=
y(e^x)^2-y.
\intertext{Umgeschrieben als quadratische Gleichung für $e^x$ ist dies}
(1-y) (e^x)^2
&=
1 +y
&
&&
(1-y) (e^x)^2
&=
-1-y.
\intertext{Mit der Quadratwurzel findet man}
e^x
&=
\!\sqrt{\frac{1+y}{1-y}}
&
&\text{bzw.}&
e^x
&=
\!\sqrt{\frac{y+1}{y-1}}
\intertext{und schliesslich mit dem Logarithmus den Wert}
x
=
\operatorname{artanh}y
&=
\log
\!\sqrt{\frac{1+y}{1-y}}
&
&\text{bzw.}&
x
=
\operatorname{arcoth}y
&=
\log
\!\sqrt{\frac{y+1}{y-1}}
\\
&=
\frac12\log
\frac{1+y}{1-y}
&
&&
&=
\frac12\log
\frac{y+1}{y-1}
\end{align*}
der Umkehrfunktion der hyperbolischen Tangens- bzw.~Kotangensfunktion.
\index{artanh@$\operatorname{artanh}$}%
\index{arcoth@$\operatorname{arcoth}$}%

\begin{satz}[hyperbolische Umkehrfunktionen]
\label{buch:intalg:elementar:trigo:satz:invhyperbolisch}
Die Umkehrfunktionen der hyperbolischen Funktionen sind
\index{hyperbolische Funktion!Umkehrfunktion}%
\begin{align*}
\operatorname{arsinh}y
&=
\log\Bigl(
y + \!\sqrt{\smash{y^2}+1\raisebox{1.5pt}{\mathstrut}}
\Bigr),
&
\operatorname{arcosh}y
&=
\log\Bigl(
y + \!\sqrt{\smash{y^2}-1\raisebox{1pt}{\mathstrut}}
\Bigr),
\\
\operatorname{artanh}y
&=
\frac12\log
\frac{1+y}{1-y},
&
\operatorname{arcoth}y
&=
\frac12\log
\frac{y+1}{y-1}.
\end{align*}
\index{arsinh@$\operatorname{arsinh}$}%
\index{arcosh@$\operatorname{arcosh}$}%
\index{artanh@$\operatorname{artanh}$}%
\index{arcoth@$\operatorname{arcoth}$}%
\end{satz}

Der Satz zeigt also, dass nicht nur die hyperbolischen Funktionen
sondern auch ihre Umkehrfunktionen ausschliesslich durch algebraische
Funktionen und den Logarithmus definiert werden können.

Die Funktionen $\operatorname{tanh}x$ und $\operatorname{coth}x$ sind
reziprok.
Die beiden Funktionen $\operatorname{artanh}y$ und
$\operatorname{arcoth}y$
müssen also unter der Substitution $y\to 1/y$ ineinander übergehen.
Tatsächlich wird das Argument des Logarithmus in der Areatangensfunktion
unter der Substitution zu
\[
\frac{1+y}{1-y}
\to
\frac{1+1/y}{1-1/y}
=
\frac{y+1}{y-1},
\]
dem Argument des Logarithmus der Areakotangensfunktion.

%
% Sinus und Kosinus
%
\subsubsection{Sinus und Kosinus}
Aus der Identität
\[
e^{ix}
=
\cos x + i \sin x
\]
folgt zunächst mittels der Substitution $x\to -x$
\[
e^{-ix}
=
\cos x - i \sin x.
\]
Diese beiden Formeln bilden ein lineares Gleichungssystem
\index{lineares Gleichungssystem}%
\[
\renewcommand{\arraycolsep}{2pt}
\begin{array}{rcrcl}
\cos x &+& i\sin x &=& e^{\phantom{-}ix} \\
\cos x &-& i\sin x &=& e^{-ix}
\end{array}
\]
für die Funktionen $\cos x$ und $\sin x$.
\index{sin@$\sin$}%
\index{cos@$\cos$}%
Es kann durch Addition bzw.~Subtraktion gelöst werden.
Die Lösungen sind
\[
\cos x = \frac{e^{ix}+e^{-ix}}{2}
\qquad\text{und}\qquad
\sin x = \frac{e^{ix}-e^{-ix}}{2i}.
\]
Die Funktionen können also mit der Exponentialfunktion in der gleichen
Form ausgedrückt werden, in der die hyperbolischen Funktionen 
definiert sind.
Durch Vergleich mit den Formeln in der
Definition~\ref{buch:intalg:elementar:trigo:def:hyperbolisch}
ergeben sich die Beziehungen
\[
\cos x = \cosh ix
\qquad\text{und}\qquad
\sin x = \frac{1}{i} \sinh ix.
\]
Die Multiplikation mit $i$ entspricht geometrisch einer Drehung
der komplexen Zahlenebene um $90^\circ$.
Man kann also sagen, dass trigonometrischen Funktionen durch Drehung
der Argumentebene um $90^\circ$ aus den hyperbolischen Funktionen
hervorgehen.

Für die Umkehrfunktionen $\arccos y$ und $\arcsin y$ sind die gleichen
algebraischen Umformungen anwendbar, die auf die Umkehrfunktionen
$\operatorname{arcosh}y$ und $\operatorname{arsinh}y$ der hyperbolischen
Funktionen geführt haben.
\index{arccos@$\arccos$}%
\index{arcsin@$\arcsin$}%
Aus
\input{chapters/070-intalg/fig/fig-invcis.tex}%
\begin{align}
\cos x = \frac{e^{ix}+e^{-ix}}{2} &= y
&
&\text{bzw.}&
\sin x = \frac{e^{ix}-e^{-ix}}{2i} &= y
\notag
\intertext{folgen durch Multiplikation mit dem Nenner und mit $e^{ix}$
die quadratischen Gleichungen}
(e^{ix})^2 -2ye^{ix} + 1 &= 0
&
&&
(e^{ix})^2 - 2iy e^{ix} - 1 &= 0
\notag
\intertext{für $e^{ix}$ mit den Lösungen}
e^{ix}
&=
y \pm \!\sqrt{\smash{y^2}-1\raisebox{1.5pt}{\mathstrut}}
&
&&
e^{ix}
&=
iy\pm \!\sqrt{-\smash{y^2}+1\raisebox{1.5pt}{\mathstrut}}.
\notag
\intertext{Da die Kosinuswerte Betrag nicht grösser als $1$ haben
können, ist der Radikand in der linken Gleichung immer negativ.
Wir können das Vorzeichen kehren, indem wir den Faktor $i$ vor die
Wurzel ziehen:}
e^{ix}
&=
y \pm i\!\sqrt{1-\smash{y^2}\raisebox{1.5pt}{\mathstrut}}
&
&&
e^{ix}
&=
iy\pm \!\sqrt{1-\smash{y^2}\raisebox{1.5pt}{\mathstrut}}.
\label{buch:intalg:elementar:trigo:eqn:invcisrhs}
\end{align}
Abbildung~\ref{buch:intalg:elementar:trigo:fig:invcis} zeigt 
in rot die rechten Seiten der Gleichungen
\eqref{buch:intalg:elementar:trigo:eqn:invcisrhs}.
Da die Punkte auf dem Einheitskreis liegen, ist der Logarithmus
der in blau eingezeichnete Winkel $x$.

\begin{satz}
\label{buch:intalg:elementar:trigo:satz:trigonometrisch}
Die trigonometrischen Funktionen
\begin{align*}
\cos x &= \frac{e^{ix}+e^{-ix}}{2}
&&\text{und}&
\sin x &= \frac{e^{ix}-e^{-ix}}{2i}
\intertext{in Exponentialdarstellung haben die logarithmischen
Umkehrfunktionen}
\arccos y
&=
\frac{1}{i}
\log
\Bigl(
y + i\!\sqrt{1-\smash{y^2}\raisebox{1.5pt}{\mathstrut}}
\Bigr)
&&\text{bzw.}&
\arcsin y
&=
\frac{1}{i}
\log
\Bigl(
iy + \!\sqrt{1-\smash{y^2}\raisebox{1.5pt}{\mathstrut}}
\Bigr).
\end{align*}
\end{satz}

%
% Tangens und Kotangens
%
\subsubsection{Tangens und Kotangens}
Tangens und Kotangens sind Quotienten der Sinus- und Kosinusfunktionen
\input{chapters/070-intalg/fig/fig-invtan.tex}%
oder in Exponentialdarstellung
\index{tan@$\tan$}%
\index{cot@$\cot$}%
\begin{align*}
\tan x
&=
\frac1{i}\frac{e^{ix}-e^{-ix}}{e^{ix}+e^{-ix}}
&
&\text{bzw.}&
\cot x
&=
i\frac{e^{ix}+e^{-ix}}{e^{ix}-e^{-ix}}
\end{align*}
Für die Umkehrfunktionen setzt man wieder 
\index{arctan}%
\begin{align*}
\tan x
= 
\frac1i\frac{e^{ix}-e^{-ix}}{e^{ix}+e^{-ix}}
&=
y
&&&
\cot x
=
i\frac{e^{ix}+e^{-ix}}{e^{ix}-e^{-ix}}
&=
y
\intertext{und multipliziert mit dem Nenner}
e^{ix}-e^{-ix}
&=
iy
(e^{ix}+e^{-ix})
&&&
e^{ix}+e^{-ix}
&=
-iy
(e^{ix}-e^{-ix})
\\
(1-iy)e^{ix}
&=
(1+iy)e^{-ix}
&&&
(iy-1)e^{ix}
&=
-(1+iy)e^{-ix}.
\intertext{Durch Multiplizieren mit $e^{ix}$ kann man nach $e^{2ix}$
auflösen und findet}
e^{2ix}
&=
\frac{1+iy}{1-iy}
&&&
e^{2ix}
&=
\frac{1-iy}{1+iy}
\\
x
&=
\frac1{2i}
\log
\frac{1+iy}{1-iy}
&&\text{bzw.}&
x
&=
\frac1{2i}
\log
\frac{1-iy}{1+iy}.
\end{align*}
In Abbildung~\ref{buch:intalg:elementar:trigo:fig:invtan} sind
die Zahlen $1+iy$ und $1-iy$ dargestellt.
Der gesuchte Winkel $x$ tritt beim Ursprung auf.
Da die Multiplikation mit einer komplexen Zahl eine Drehungstreckung ist,
muss der Quotient den Drehstreckungsfaktor zurückliefen.
Der Quotient von $1+iy$ und $1-iy$ ist ein komplexe Zahl vom Betrag $1$,
die Drehung von $1-iy$ auf $i+iy$ beschreibt.
Der zugehörige Drehwinkel ist $2x$.
Somit ist der gesuchte Winkel tatsächlich der halbe komplexe Logarithmus
des Quotienten der beiden Zahlen.

\begin{satz}
\label{buch:intalg:elementar:satz:trigo:tancot}
Die trigonometrischen Funktionen
\begin{align*}
\tan x
&=
\frac1{i}\frac{e^{ix}-e^{-ix}}{e^{ix}+e^{-ix}}
&
&\text{und}&
\cot x
&=
i\frac{e^{ix}+e^{-ix}}{e^{ix}-e^{-ix}}
\intertext{in Exponentialdarstellung haben die logarithmische Umkehrfunktion}
\arctan y
&=
\frac1{2i}
\log
\frac{1+iy}{1-iy}
&&\text{bzw.}&
\operatorname{arccot} y
&=
\frac1{2i}
\log
\frac{1-iy}{1+iy}.
\end{align*}
\end{satz}

%
% Stammfunktion einer rationalen Funktion
%
\subsubsection{Stammfunktion einer rationalen Funktion}
In Abschnitt~\ref{buch:integration:residuen:section:beispiel}
wurde die Aufgabe gestellt, das die Stammfunktion
\[
\int \frac{dx}{x^2+1}
\]
zu bestimmen.
Die Partialbruchzerlegung des Integranden ist
\[
\frac{1}{x^2+1}
=
\frac1{2i}
\biggl(
\frac{1}{x-i}
-
\frac{1}{x+i}
\biggr).
\]
In dieser Form ist die Stammfunktion einfach zu bestimmen, sie ist
\begin{align*}
\int
\frac{1}{x^2+1}
\,dx
&=
\frac{1}{2i}
\bigl(
\log(x-i)
-
\log(x+i)
\bigr)
=
\frac{1}{2i}
\log
\frac{x-i}{x+i}
=
\frac{1}{2i}
\log
\frac{1+ix}{-1+ix}.
\intertext{%
In dieser Form stimmt das Resultat nicht mit der Form des erwarteten
Resultates, nämlich der Arkustangens-Funktion, überein.
Der Nenner hat das falsche Vorzeichen.
Multiplikation mit $-1$ im Argument des Logarithmus ändert aber den
Wert der Logarithmusfunktion nur um die Konstante $\pi$, also kann
man immer noch schliessen}
&=
\frac{1}{2i}
\log
\frac{1+ix}{1-ix}
+C
=
\arctan x + C.
\end{align*}
Das Beispiel illustriert, dass sich Integrale von rationalen Funktionen
über $\mathbb{R}$ tatsächlich ausschliesslich mit der komplexen
Logarithmus-Funktion berechnen lassen.

%
% Die Differentialkörper der hyperbolischen Funktionen
%
\subsubsection{Die Differentialkörper der hyperbolischen Funktionen}
Der Körper $\mathbb{Q}(x)$ enthält weder die trigonometrischen noch
die hyperbolischen Funktionen.
Wie früher betrachten wir die Derivation $D$ mit $Dx=1$, die die
Ableitung nach $x$ symbolisiert.

Um die hyperbolischen Funktionen zu $\mathbb{Q}(x)$ hinzuzufügen,
muss ein Exponentialfunktion $e(x)$ hinzugefügt werden.
Damit lassen sich dann die Funktionen
\[
\cosh x
=
\frac12
\biggl(
e(x) + \frac{1}{e(x)}
\biggr),
\quad
\sinh x
=
\frac12
\biggl(
e(x) - \frac{1}{e(x)}
\biggr),
\quad
\tanh x
=
\frac{
\displaystyle
e(x)-\frac{1}{e(x)}
}{
\displaystyle
e(x)+\frac{1}{e(x)}
}
=
\frac{e(x)^2-1}{e(x)^2+1}
.
\]
Die bekannte Relation
\[
(\cosh x)^2
-
(\sinh x)^2
=
\frac14\biggl(
e(x)^2 + 2e(x)\frac{1}{e(x)} + \frac{1}{(e(x))^2}
\biggr)
-
\frac14\biggl(
e(x)^2 - 2e(x)\frac{1}{e(x)} + \frac{1}{(e(x))^2}
\biggr)
=
1
\]
folgt allein durch algebraische Umformung.
Die Ableitungseigenschaften der Funktion $e(x)$ im Differentialkörper
als Exponentialfunktion von $x$ reichen aus, um die Ableitungsregeln
\begin{align*}
D\cosh x
&=
\frac12\biggl(
De(x) + D\frac1{e(x)}
\biggr)
=
\frac12
\biggl(
e(x) - \frac{1}{e(x)}
\biggr)
=
\sinh x,
\\
D\sinh x
&=
\frac12
\biggl(
De(x) - D\frac{1}{e(x)}
\biggr)
=
\frac12
\biggl(
e(x) + \frac{1}{e(x)}
\biggr)
=
\cosh x
\end{align*}
für die hyperbolischen Funktionen zu verifizieren.
Auch die Ableitungsregel für den hyperbolischen Tangens können wir
leicht nachrechnen:
\begin{align*}
D\tanh x
&=
D
\frac{\sinh x}{\cosh x}
\\
&=
\frac{(D\sinh x)\cdot \cosh x - \sinh x\cdot(D\cosh x)}{(\cosh x)^2}
\\
&=
\frac{(\cosh x)^2 - (\sinh x)^2}{(\cosh x)^2}
\\
&=
\frac{1}{(\cosh x)^2}.
\end{align*}
Dies ist die bekannte Ableitung der hyperbolischen Tangensfunktion.
Diese Beispiele zeigen, dass alle Eigenschaften der hyperbolischen
Funktionen von den Details der Funktion $e(x)$ unabhängig sind.

Für die hyperbolischen Umkehrfunktionen müssen gemäss den Formeln
von Satz~\ref{buch:intalg:elementar:trigo:satz:invhyperbolisch}
zunächst die Wurzeln
\[
s_+(x) = \sqrt{\smash{y^2} + 1\raisebox{1.5pt}{\mathstrut}}
\qquad\text{und}\qquad
s_-(x) = \sqrt{\smash{y^2} - 1\raisebox{1.5pt}{\mathstrut}}
\]
hinzufügen, mit denen anschliessend im Funktionenkörper
$\mathbb{Q}(x,e(x),s_+(x),s_-(x))$
die Funktionen
\[
x + s_+(x)
\qquad\text{und}\qquad
x + s_-(x)
\]
gebildet werden können.
Indem wir die Logarithmen $\operatorname{arsinh}(x) = \log(x+s_+(x))$
und $\operatorname{arcosh}(x) = \log(x+s_-(x))$
hinzufügen, erhalten wir den Differentialkörper
\[
\mathbb{Q}(x,e(x),s_\pm(x),
\operatorname{arsinh}(x),\operatorname{arcosh}(x)),
\]
die die gesuchten Umkehrfunktionen entält.
Die hyperbolischen Funktionen und ihre Umkehrfunktionen können also
allein durch den Prozess der algebraischen, exponentiellen oder
logarithmischen Erweiterung erhalten werden.

%
% Die Differentialkörper der trigonometrischen Funktionen
%
\subsubsection{Die Differentialkörper der trigonometrischen Funktionen}
Ein Differentialkörper für die trigonometrischen Funktionen und ihre
Umkehrfunktionen lässt sich ganz analog zur Konstruktion für die 
hyperbolischen Funktionen im vorangegangenen Abschnitt konstruieren.
Damit die Konstruktion möglich wird, muss dem Körper $\mathbb{Q}(x)$
im ersten Schritt die neue Konstante $i$ durch algebraische
Erweiterung mit dem Minimalpolynom $x^2+1$ hinzugefügt werden.
Der Ausgangskörper für die Konstruktion ist also $\mathbb{Q}(x,i)$.

Statt der reellen Exponentialfunktion fügen wir im nächsten Schritt 
dem Körper die Exponentialfunktion $e(x)$ von $ix$ hinzu.
Sie erfüllt die definierende Eigenschaft $De(x) = e(x)\cdot D(ix) = ie(x)$.
Damit können dann die trigonometrischen Funktionen
\begin{align*}
\sin x
&=
\frac12\biggl(
e(x)+\frac{1}{e(x)}
\biggr),
&
\cos x
&=
\frac1{2i}\biggl(
e(x)-\frac{1}{e(x)}
\biggr),
&
\tan x
&=
\frac{1}{i}
\frac{
\displaystyle e(x)-\frac{1}{e(x)}
}{
\displaystyle e(x)+\frac{1}{e(x)}
}
=
\frac{1}{i}\frac{e(x)^2-1}{e(x)^2+1}
\end{align*}
durch algebraische Operationen in diesem Körper konstruiert werden.
Der ``Satz des Pythagoras'' ist eine rein algebraische Eigenschaft,
es gilt nämlich
\begin{align*}
(\sin x)^2 + (\cos x)^2
&=
\frac1{4i^2}
\biggl(
e(x)^2 - 2 + \frac{1}{e(x)^2}
\biggr)
+
\frac14
\biggl(
e(x)^2 + 2 + \frac{1}{e(x)^2}
\biggr)
\\
&=
-
\frac1{4}
\biggl(
e(x)^2 - 2 + \frac{1}{e(x)^2}
\biggr)
+
\frac14
\biggl(
e(x)^2 + 2 + \frac{1}{e(x)^2}
\biggr)
=
1.
\end{align*}
Auch die Ableitungsregeln folgen direkt aus der Definition
\begin{align*}
D\sin x
&=
\frac{1}{2i}
\biggl(
De(x) - D\frac{1}{e(x)}
\biggr)
=
\frac{1}{2i}
\biggl(
ie(x) + i\frac{1}{e(x)}
\biggr)
=
\cos x
\\
D\cos x
&=
\frac12
\biggl(
De(x) + D\frac{1}{e(x)}
\biggr)
=
\frac12
\biggl(
ie(x) - i\frac{1}{e(x)}
\biggr)
=
-
\frac1{2i}
\biggl(
e(x) - \frac{1}{e(x)}
\biggr)
=
-
\sin x.
\end{align*}
Die Funktionen $\sin x$ und $\cos x$ sind also in der Körpererweiterung
$\mathbb{Q}(x,i,e(x))$ enthalten.

Für die Umkehrfunktionen ist nach
Satz~\ref{buch:intalg:elementar:trigo:satz:trigonometrisch}
zusätzlich das algebraische Element
\[
s(x)
=
\! \sqrt{1-\smash{x^2}\raisebox{1.5pt}{\mathstrut}}
\]
hinzuzufügen.
Der Differentialkörper $\mathbb{Q}(x,i,s(x))$ enthält dann auch
die Funktionen
\[
x+is(x)
=
x + i\!\sqrt{1-\smash{y^2}\raisebox{1.5pt}{\mathstrut}}
\qquad
\text{und}
\qquad
ix+s(x)
=
ix + \!\sqrt{1-\smash{y^2}\raisebox{1.5pt}{\mathstrut}},
\]
da sie durch algebraische Operationen entstehen.
Schliesslich müssen nur noch die Logarithmen
\[
\vartheta_1(x)
=
\log(y+is(x))
\qquad
\text{und}
\qquad
\vartheta_2(x)
=
\log(iy+s(x))
\]
dieser Funktionen hinzugefügt werden, damit dann in
$\mathbb{Q}(x,i,s(x),\vartheta_1(x),\vartheta_2(x))$
die Funktionen
\[
\arccos x = \frac{1}{i}\vartheta_1(x)
\qquad
\text{und}
\qquad
\arcsin x = \frac{1}{i}\vartheta_2(x)
\]
enthält.

Wie bei den hyperbolischen Funktionen können wir schliessen, dass
der Differentialkörper $\mathbb{Q}(x,i,s(x),\vartheta_1(x),\vartheta_2(x))$
die inversen trigonometrischen Funktionen enthält.
Der hier vorgestellte rein algebraische Rahmen enthält also alle
Funktionen, denen man im Rahmen einer Analysis-Vorlesung begegnet.
%Die trigonometrischen Funktionen ebenso wie ihre Umkehrfunktionen
%sind daher elementare Funktionen.

%
% Eindeutigkeit der trigonometrischen und hyperbolischen Funktionen
%
\subsubsection{Eindeutigkeit der trigonometrischen und hyperbolischen
Funktionen}
Die Definition der trigonometrischen und hyperbolischen Funktionen 
sowie ihrer Umkehrfunktionen in einem Differentialkörper basieren
auf einer Exponential- und Logarithmus-Funktionen.
Da diese nicht eindeutig bestimmt sind, sind auch die
trigonometrischen und hyperbolischen Funktionen nicht eindeutig
bestimmt.
Die Exponentialfunktion ist nur bis auf einen konstanten Faktor bestimmt,
der sich natürlich auch in der Definition der trigonometrischen und
hyperbolischen Funktionen niederschlägt.
Dies scheint den Beziehungen
\[
\cos^2 x + \sin^2 x = 1
\qquad\text{und}\qquad
\cosh^2 x - \sinh^2 x = 1
\]
zu widersprechen, die für jede Wahl der Exponentialfunktion gelten.
Wir versuchen daher für die konkrete Wahl der Exponentialfunktion $e^x$
der klassischen Analysis den Einfluss eines Faktors auf die
trigonometrischen und hyperbolischen Funktionen zu ermitteln.
Im Folgenden reservieren wir daher die Notation $e^f$ und $\log f$
für die klassische definierten Exponential- und Logarithmusfunktionen.
Ebenso unterscheiden wir die verallgemeinerten trigonometrischen
und hyperbolischen Funktion mithilfe eines grossen Anfangsbuchstabens
von den klassischen Funtkionen.

Wir beginnen mit den hyperbolischen Funktionen.
Eine beliebige Exponentialfunktion $e(x)$ von $x$ ist ein Vielfaches
$a e^x$ der klassischen Exponentialfunktion.
Die hyperbolischen Funktionen sind dann
\begin{align*}
\operatorname{Cosh} x
&=
\frac{ae^x + e^{-x}/a}{2}
&&\text{und}&
\operatorname{Sinh} x
&=
\frac{ae^x - e^{-x}/a}{2}.
\intertext{Mit den Exponentialgesetzen folgt}
\operatorname{Cosh} x
&=
\frac{e^{x+\log a} + e^{-x-\log a}}{2}
&& \text{und}&
\operatorname{Sinh} x
&=
\frac{e^{x+\log a} - e^{-x-\log a}}{2}
\\
&=
\cosh(x+\log a)
&&&
&=
\sinh(x+\log a).
\end{align*}
Die verschiedenen möglichen hyperbolischen Funktionen entstehen 
durch Verschiebung der unabhängigen Variablen.
Da man die Funktionen $\operatorname{Cosh}$ und $\operatorname{Sinh}$
\index{Cosh@$\operatorname{Cosh}$}
\index{Sinh@$\operatorname{Sinh}$}
\index{autonom}%
durch eine autonome\footnote{Eine autonome Differentialgleichung ist eine Differentialgleichung der Form $y'=f(y)$, die nicht explizit von $x$ abhängt.
Sie zeichnen sich dadurch aus, dass mit $y(x)$ auch die verschobene
Funktion $y(x+\delta)$ eine Lösung ist.} Differentialgleichung definiert
ansehen kann, ist nachvollziehbar, dass die Definition die verschiedenen
möglichen Lösungen nicht unterscheiden kann.

Bei den trigonometrischen Funktionen betrachten wir einen Faktor der
Form $a=e^{i\delta}$, woraus $1/a = e^{-i\delta}$ folgt.
Die trigonometrischen Funktionen sind dann
\index{Cos@$\operatorname{Cos}$}
\index{Sin@$\operatorname{Sin}$}
\begin{align*}
\operatorname{Cos} x
&=
\frac12(ae^{ix}+e^{-ix}/a)
=
\frac12(e^{i\delta}e^{ix}+e^{ix}e^{-i\delta})
=
\frac12(e^{i(x+\delta)}+e^{-i(x+\delta)})
=
\cos(x+\delta)
\\
\operatorname{Sin} x
&=
\frac12(ae^{ix}-e^{-ix}/a)
=
\frac12(e^{i\delta}e^{ix}-e^{ix}e^{-i\delta})
=
\frac12(e^{i(x+\delta)}-e^{-i(x+\delta)})
=
\sin(x+\delta).
\end{align*}
Multiplikation mit einem Skalenfaktor bedeutet also wieder nur eine
Verschiebung des Arguments.

Die oben durchgeführten Rechnungen funktionieren natürlich auch für
einen komplexen Faktor, den wir $a=e^{b}$ schreiben können.
\begin{align*}
\operatorname{Cosh} x
&=
\frac12 ( ae^{x} + e^{-x}/a )
=
\frac12 ( e^{x+b} + e^{-x-b} )
=
\cosh(x+b)
\\
\operatorname{Sinh} x
&=
\frac12 ( ae^{x} - e^{-x}/a )
=
\frac12 ( e^{x+b} - e^{-x-b} )
=
\sinh(x+b).
\end{align*}
Entsprechende Formeln für die trigonometrischen Funktionen ergeben sich
daraus durch Multiplikation des Argumentes mit $i$.
In allen Fällen führen die verschiedenen Expontialfunktion zu einer
möglicherweise um einen komplexen Summanden verschobenen Funktion.

Bei den Umkehrfunktionen ist die Situation noch einfacher.
Sie sind als ein Logarithmus definiert.
Ein solcher ist nur bis auf einen konstanten Summanden definiert.
Diese Beobachtung ist auch konsistent mit der Interpretation der
Umkehrbeziehung für die Exponential- und Logarithmusfunktionen.
Die additive Konstante, um die sich die trigonometrischen und
hyperbolischen Umkehrfunktionen unterscheiden, ist auch die konstante
Translation, um die sich die trigonometrischen und hyperbolischen
Funktionen unterscheiden.
